\section{API et bases de données}
\label{sec:backend}

Cette section décrit les choix techniques retenus pour le backend de l'application : l'architecture mise en place, la structure de la base de données et le fonctionnement de l'API.

\subsection{Architecture générale du backend}

Le backend est développé avec \textbf{FastAPI}, un framework Python conçu pour la création d'API REST. Pour la persistance des données, nous avons choisi \textbf{PostgreSQL}, piloté via l'ORM \textbf{SQLAlchemy} qui permet de manipuler la base de données sous forme d'objets Python, ce qui simplifie la lecture et la maintenance du code. L'ensemble de l'infrastructure est conteneurisé avec \textbf{Docker}, ce qui garantit un environnement identique entre le développement et la production.

La structure du projet suit une architecture modulaire en couches :
\begin{itemize}
    \item \texttt{models/} : Définition des entités et des tables via SQLAlchemy.
    \item \texttt{schemas/} : Validation et typage des données entrantes et sortantes via Pydantic.
    \item \texttt{routers/} : Définition des points d'accès (endpoints) de l'API.
    \item \texttt{utils/} : Fonctions transversales partagées, notamment la sécurité et le hachage.
\end{itemize}

\subsection{Base de données}

La base de données centralise les données des utilisateurs et assure la persistance de leur activité de navigation.

\subsubsection{Modélisation des données}

Nous avons modélisé les entités autour des besoins concrets de l'application : les comptes utilisateurs, leurs vélos, les itinéraires calculés, l'historique de navigation et les signalements communautaires. Chaque entité correspond à une table PostgreSQL, créée automatiquement au démarrage du serveur via SQLAlchemy.

\subsubsection{Schéma relationnel}

Les entités sont liées entre elles par des relations qui reflètent la logique métier de l'application :
\begin{itemize}
    \item Un \textbf{utilisateur} peut posséder plusieurs \textbf{vélos} (relation 1-N).
    \item Un \textbf{utilisateur} peut sauvegarder plusieurs \textbf{itinéraires} (relation 1-N).
    \item Un \textbf{itinéraire} peut être associé à plusieurs entrées dans l'\textbf{historique} de navigation (relation 1-N).
    \item Un \textbf{utilisateur} peut émettre plusieurs \textbf{signalements} de dangers (relation 1-N).
\end{itemize}

\subsubsection{Description des tables}

Voici le rôle de chaque table dans la base de données :

\begin{itemize}
    \item \textbf{Table \texttt{users}} : Stocke les données de profil, dont les identifiants de connexion hachés et les informations personnelles.
    \item \textbf{Table \texttt{bikes}} : Contient les caractéristiques techniques des vélos (type, assistance électrique) rattachés à chaque utilisateur.
    \item \textbf{Table \texttt{routes}} : Enregistre les trajets calculés avec leurs coordonnées, distance, durée estimée et indice de sécurité.
    \item \textbf{Table \texttt{user\_history}} : Retrace l'activité de navigation en associant chaque trajet effectué à un horodatage.
    \item \textbf{Table \texttt{reports}} : Conserve les signalements géolocalisés émis par la communauté (accidents, travaux, obstacles).
    \item \textbf{Table \texttt{road\_accidents}} : Recense les accidents corporels géolocalisés publiés par les sources officielles (BAAC en France, Statbel en Belgique), à ne pas confondre avec \texttt{reports} : les premiers sont des faits historiques établis, les seconds des signalements ponctuels d'utilisateurs. Ils corrigent la note de sécurité des segments concernés (cf. Section~\ref{sec:accidentologie}).
    \item \textbf{Tables \texttt{badges} et \texttt{user\_badges}} : Posent les bases d'un système de gamification prévu pour une version ultérieure.
\end{itemize}

\subsection{Gestion du trafic en temps réel}

La sécurité d'un itinéraire cyclable dépend en partie des conditions de circulation au moment du trajet. Un axe embouteillé concentre des manœuvres imprévisibles et augmente le risque d'accident pour le cycliste. Notre système intègre donc une couche de données trafic en temps réel, qui alimente à la fois l'affichage sur la carte et le calcul d'itinéraire.

\subsubsection{Choix de la source : une comparaison mesurée}

Deux sources ouvertes couvrent la métropole bordelaise : l'\textbf{API AVATAR du Cerema}, qui expose les mesures brutes de points de comptage fixes (boucles électromagnétiques et radars), et l'\textbf{open data de Bordeaux Métropole} (jeu \texttt{ci\_trafi\_l}), qui publie directement l'état de congestion de chaque tronçon de voirie. Nos deux chaînes de traitement s'étaient initialement appuyées sur l'une et l'autre sans concertation ; nous les avons départagées en instrumentant les deux API.

\begin{itemize}
    \item \textbf{Latence de réponse} : 0,45\,s pour Bordeaux Métropole, contre 31\,s pour AVATAR, dont l'API n'offre aucun filtre spatial et impose donc de télécharger les 9\,754 points de comptage nationaux.
    \item \textbf{Fraîcheur} : Bordeaux Métropole publie toutes les 5 minutes exactement, l'horodatage correspondant à l'instant de publication. Les mesures AVATAR accusent 10 à 15 minutes de retard en régime normal, et nous avons observé des interruptions de flux dépassant l'heure.
    \item \textbf{Complétude} : sur les 602 capteurs AVATAR de la zone, 403 seulement remontent des mesures ; le taux d'occupation n'est jamais renseigné et la vitesse ne l'est que dans 21\,\% des cas — or c'était précisément la variable la plus pondérée par notre indice de danger.
    \item \textbf{Nature de la donnée} : Bordeaux Métropole fournit un état déjà qualifié, porté par la géométrie du tronçon, là où AVATAR impose de reconstruire un score à partir de mesures ponctuelles.
\end{itemize}

Nous avons donc écarté les comptages bruts au profit de l'\textbf{état de congestion déjà qualifié sur une géométrie de tronçon}, dont Bordeaux Métropole a servi de référence. La couverture nationale d'AVATAR, seul argument en sa faveur, s'est révélée trompeuse : elle ignore Toulouse, Lille et Nice, et ne dépasse pas les frontières, ce qui la rend inutile pour notre profil belge de Tournai.

\subsubsection{Généralisation à plusieurs métropoles}

Ce format n'est pas propre à Bordeaux : plusieurs métropoles françaises publient, sur la même plateforme \textit{Opendatasoft} (API \texttt{explore/v2.1}), un jeu de tronçons routiers assortis d'un état de congestion pré-qualifié. Nous avons donc généralisé la collecte à un \textbf{registre de sources} : chaque entrée décrit une métropole par l'URL de son jeu, son emprise géographique, et la correspondance entre son vocabulaire d'état et notre échelle unique (\textit{fluide} / \textit{dense} / \textit{embouteillé} / \textit{indéterminé}). Ajouter une ville se réduit alors à une entrée de configuration.

Quatre sources sont ainsi intégrées : \textbf{Bordeaux Métropole} (état \texttt{FLUIDE}/\texttt{DENSE}/\texttt{EMBOUTEILLE}), \textbf{Rennes Métropole} (champ \texttt{trafficstatus} au standard européen DATEX~II), \textbf{Eurométropole de Strasbourg} (système SIRAC, état codé de 0 à 3) et \textbf{Nantes Métropole} (indicateur de fluidité des axes). Un unique collecteur générique les interroge toutes, quel que soit leur vocabulaire d'origine. À l'inverse, nous avons écarté les sources qui ne publient que des comptages bruts (Paris, Lille, Toulouse), celle de Lyon dont le format n'est pas interopérable avec les précédentes, et l'ensemble des sources belges, faute d'état qualifié sur la voirie urbaine.

\subsubsection{Une collecte, deux usages}

La collecte est menée par une tâche de fond, au rythme de publication de la source, et jamais dans le fil d'une requête utilisateur : l'endpoint \texttt{/traffic/} se contente de servir l'instantané déjà constitué, ce qui rend sa réponse immédiate même au premier appel. Le même instantané produit deux sorties : la couche GeoJSON affichée sur la carte, et le drapeau de congestion posé sur les arêtes du graphe. L'utilisateur voit ainsi exactement ce que le calcul d'itinéraire cherche à éviter.

Une source n'est interrogée que si l'emprise qu'elle couvre croise celle du graphe chargé, elle-même déduite des nœuds du graphe. Le critère est délibérément géographique et non national : notre profil « tournai » comprend des communes françaises (Lille, Roubaix), et un filtrage par pays y aurait activé la source bordelaise pour ne récolter aucun tronçon atteignant la zone. Sur un profil non couvert, la couche se désactive proprement, l'option disparaît de l'interface et aucune requête n'est émise.

\subsubsection{Rattachement au graphe}

Les tronçons signalés \textit{dense} ou \textit{embouteillé} sont échantillonnés tous les 20 mètres, et chaque point est rattaché au nœud du graphe le plus proche au moyen de l'index spatial (\textit{BallTree}) déjà précalculé pour le routage. Une arête n'est déclarée congestionnée que si \textbf{ses deux extrémités} figurent parmi ces nœuds : se contenter d'une seule extrémité marquait toutes les rues perpendiculaires débouchant sur l'axe, soit près de trois fois plus d'arêtes que nécessaire.

En sont exclues les voies \textbf{physiquement séparées de la chaussée} (pistes cyclables en site propre, chemins) : le cycliste n'y partage pas la circulation automobile, si bien qu'une piste longeant un axe saturé ne doit pas en hériter la pénalité.

\subsubsection{Malus de routage}

Le malus est \textbf{gradué selon la sévérité} : un axe embouteillé (rouge) pèse plein tarif, un axe dense (orange) 40\,\% seulement. Il se décompose en une part fixe, qui s'applique quel que soit le curseur rapidité/sécurité car un embouteillage ralentit aussi le trajet, et une part qui ne croît qu'en allant vers la sécurité. Concrètement, un tronçon rouge coûte de 50 mètres virtuels supplémentaires (rapidité pure) à 170 (sécurité pure), un tronçon orange de 20 à 68.

Ce dimensionnement est volontairement modéré. Les axes congestionnés sont massivement des routes primaires et secondaires, que le score d'infrastructure évite déjà fortement — de l'ordre de 700 mètres virtuels sur une arête, en mode sécurité. Le trafic n'ajoute qu'environ un quart par-dessus : il affine le choix sans constituer un second évitement des grands axes, qui ferait double emploi. Nous avons par ailleurs renoncé à faire du niveau de congestion un proxy du risque d'accident : pour un cycliste, un axe à l'arrêt n'est pas plus dangereux qu'un axe dense mais roulant, la gravité ne suivant pas l'intensité de la circulation.

Enfin, le cache d'itinéraires n'est vidé que lorsque l'ensemble des arêtes congestionnées a effectivement changé, et une collecte en échec conserve le dernier état connu en le marquant comme périmé plutôt que de relâcher les pénalités.

\subsection{API REST}

L'API constitue l'interface entre les clients (web et mobile) et la logique serveur. Elle est construite autour de FastAPI et expose un ensemble d'endpoints organisés par ressource.

\subsubsection{Technologies utilisées}

FastAPI génère automatiquement une documentation interactive via Swagger UI, ce qui a facilité nos tests pendant le développement. La validation des données entrantes est assurée par \textbf{Pydantic} : chaque requête est vérifiée avant d'atteindre la logique métier, ce qui limite les erreurs en cascade.

Pour la sécurité des échanges, nous utilisons \textbf{JWT} (JSON Web Token) pour la gestion des sessions, et \textbf{Argon2} pour le hachage des mots de passe, un algorithme reconnu pour sa résistance aux attaques par force brute.

\subsubsection{Authentification et sécurité}

L'authentification repose sur le standard \textbf{OAuth2 avec Bearer Token}. À la connexion, le serveur délivre un jeton JWT que le client doit inclure dans l'en-tête de chaque requête protégée.

Nous distinguons deux niveaux d'accès :
\begin{itemize}
    \item \textbf{Accès authentifié} : obligatoire pour les actions sensibles (modification du profil, gestion des vélos, consultation de l'historique et émission de signalements).
    \item \textbf{Accès public} : le calcul d'itinéraire est accessible sans connexion, mais la sauvegarde automatique du trajet ne s'active que pour les utilisateurs connectés.
\end{itemize}

\subsubsection{Endpoints principaux}

L'API expose les endpoints suivants, organisés par ressource :

\begin{itemize}
    \item \textbf{\texttt{/users}} : Gestion des comptes, de l'inscription à la mise à jour des informations personnelles.
    \item \textbf{\texttt{/bikes}} : Gestion des vélos de l'utilisateur, avec un contrôle garantissant qu'un utilisateur ne peut accéder qu'à ses propres équipements.
    \item \textbf{\texttt{/routes}} : Calcul d'itinéraires à partir de coordonnées géographiques et des préférences de l'utilisateur (type de vélo, niveau de sécurité souhaité).
    \item \textbf{\texttt{/history}} : Consultation des trajets passés de l'utilisateur connecté.
    \item \textbf{\texttt{/stats}} : Synthèse de l'activité de l'utilisateur : distances parcourues, durée cumulée, répartition entre trajets rapides et sécurisés.
    \item \textbf{\texttt{/reports}} : Remontée de signalements par les usagers. La lecture est publique ; l'écriture requiert une authentification.
    \item \textbf{\texttt{/navigation}} : Gestion des instructions de navigation en temps réel calculer en fonction de la position actuel de l'utilisateur. En cas d'écart trop important par rapport au tracé prévu, le serveur renvoie \texttt{recalculate: true} afin que le client déclenche un nouveau calcul d'itinéraire.
\end{itemize}
